\section{Diez ejemplos de lectura}\label{sec:ejemplos}

El borde continuo delimita un grupo. El borde discontinuo identifica
notación matemática convencional. Cada leyenda de referentes es local
a su ejemplo; los índices de relaciones corresponden al catálogo de la
sección~\ref{sec:relaciones}. Las glosas explican el significado en español.

\subsection{E01. Un alcance dentro de otro}

{\small\sffamily\color{KalaAcento}Matemáticas}

\KalaMostrar{\KalaUniversal{x}{\mathbb{N}}{\KalaExistencia{y}{\mathbb{N}}{\KalaMat{x < y}}}}

\textbf{Lectura.} Para todo número natural x existe un número natural y tal que x \ensuremath{<} y.

\textbf{Interpretación.} El valor de y puede depender de x. La notación dentro de los recintos discontinuos es matemática convencional.

{\small\color{KalaSuave}No requiere leyenda de referentes.}

\vspace{6mm}

\subsection{E02. Negar una combinación}

{\small\sffamily\color{KalaAcento}Lógica y filosofía}

\KalaMostrar{\KalaNegacion{\KalaConjuncion{\KalaMat{P}}{\KalaMat{Q}}}}

\textbf{Lectura.} No es verdad que P y Q sean ambas verdaderas.

\textbf{Interpretación.} P y Q representan proposiciones. No se establece cuál de ellas es falsa.

{\small\color{KalaSuave}No requiere leyenda de referentes.}

\clearpage

\subsection{E03. Cambiar el alcance cambia el significado}

{\small\sffamily\color{KalaAcento}Lógica y filosofía}

\KalaMostrar{\KalaConjuncion{\KalaNegacion{\KalaMat{P}}}{\KalaMat{Q}}}

\textbf{Lectura.} P es falsa y Q es verdadera.

\textbf{Interpretación.} Compárala con el ejemplo anterior: la negación alcanza únicamente a P.

{\small\color{KalaSuave}No requiere leyenda de referentes.}

\vspace{6mm}

\subsection{E04. Una relación con tres argumentos}

{\small\sffamily\color{KalaAcento}Arquitectura de computadoras}

\KalaMostrar{\KalaRelacion{2}{\KalaReferencia{1}\KalaSeparador \KalaReferencia{2}\KalaSeparador \KalaReferencia{3}}}

\textbf{Lectura.} El procesador lee el dato de la dirección indicada.

\textbf{Interpretación.} Ejemplo esquemático de un único evento; no especifica cachés, coherencia ni concurrencia.

\begin{description}[font=\normalfont\bfseries,leftmargin=10mm]

\item[1] Procesador del ejemplo.

\item[2] Dirección de memoria del evento.

\item[3] Dato leído en ese evento.

\end{description}

\clearpage

\subsection{E05. Orden temporal}

{\small\sffamily\color{KalaAcento}Historia}

\KalaMostrar{\KalaRelacion{1}{\KalaReferencia{1}\KalaSeparador \KalaReferencia{2}}}

\textbf{Lectura.} El evento 1 precede temporalmente al evento 2.

\textbf{Interpretación.} Eventos ficticios en un mismo calendario. La afirmación no establece una causa.

\begin{description}[font=\normalfont\bfseries,leftmargin=10mm]

\item[1] Evento ficticio A, terminado antes de B.

\item[2] Evento ficticio B.

\end{description}

\vspace{6mm}

\subsection{E06. Partes y totalidad}

{\small\sffamily\color{KalaAcento}Biología}

\KalaMostrar{\KalaRelacion{3}{\KalaReferencia{1}\KalaSeparador \KalaReferencia{2}}}

\textbf{Lectura.} En el modelo descrito, la estructura 1 forma parte propia de la célula 2.

\textbf{Interpretación.} Ejemplo de modelado: los dos referentes son objetos concretos del mismo modelo, no clases ni conjuntos.

\begin{description}[font=\normalfont\bfseries,leftmargin=10mm]

\item[1] Una estructura identificada en el modelo.

\item[2] La célula que contiene esa estructura.

\end{description}

\clearpage

\subsection{E07. Una medición y su fuente}

{\small\sffamily\color{KalaAcento}Ciencias experimentales}

\KalaMostrar{\KalaFuente{\KalaReferencia{2}}{\KalaMat{m(S_1)=5{,}0\,\mathrm{g};\quad u(m)=0{,}1\,\mathrm{g}}}}

\textbf{Lectura.} Según el registro 2, la masa de la muestra S\textsubscript{1} es 5,0 g y su incertidumbre estándar es 0,1 g.

\textbf{Interpretación.} Datos ficticios. u(m) es incertidumbre estándar, no un error conocido ni un intervalo de confianza. S\textsubscript{1} es el identificador de muestra dentro de la fórmula.

\begin{description}[font=\normalfont\bfseries,leftmargin=10mm]

\item[2] Registro experimental ficticio; identifica la muestra S\textsubscript{1} y las condiciones de medida.

\end{description}

\vspace{6mm}

\subsection{E08. Comparar una propiedad}

{\small\sffamily\color{KalaAcento}Música y acústica}

\KalaMostrar{\KalaRelacion{4}{\KalaReferencia{1}\KalaSeparador \KalaReferencia{2}}}

\textbf{Lectura.} La frecuencia fundamental del tono 1 es mayor que la del tono 2.

\textbf{Interpretación.} Dos tonos estables en el mismo intervalo de observación. No se compara intensidad sonora.

\begin{description}[font=\normalfont\bfseries,leftmargin=10mm]

\item[1] Primer tono estable.

\item[2] Segundo tono estable.

\end{description}

\clearpage

\subsection{E09. Introducir una definición}

{\small\sffamily\color{KalaAcento}Computación y matemáticas}

\KalaMostrar{\KalaDefinicion{f}{\KalaMat{f:\mathbb{N}\to\mathbb{N};\quad f(x)=x+1}}}

\textbf{Lectura.} Se define f como la función que asigna a cada natural x su sucesor x + 1.

\textbf{Interpretación.} La definición usa notación convencional; el marcador de KalaScript identifica su función dentro del documento.

{\small\color{KalaSuave}No requiere leyenda de referentes.}

\vspace{6mm}

\subsection{E10. Una condición explícita}

{\small\sffamily\color{KalaAcento}Matemáticas}

\KalaMostrar{\KalaUniversal{n}{\mathbb{N}}{\KalaCondicional{\KalaMat{n\geq 1}}{\KalaMat{n\leq n^2}}}}

\textbf{Lectura.} Para cada número natural n, si n \ensuremath{\geq} 1, entonces n \ensuremath{\leq} n\textsuperscript{2}.

\textbf{Interpretación.} El condicional representa implicación material en este fragmento; no una orden de ejecución.

{\small\color{KalaSuave}No requiere leyenda de referentes.}
